\begin{ouverture}[statistiques]{20}{Organiser et lire des données}{Partie 5 · Modéliser et organiser}
{Tableaux · Diagrammes en barres · Courbes}
\objectif{Recueillir des données et les ranger dans un tableau.}
\objectif{Lire un tableau, un diagramme en barres ou une courbe.}
\objectif{Résoudre des problèmes à partir de données organisées.}
\prerequis{Lecture de graduations et opérations sur les entiers.}
\end{ouverture}
\pagestyle{otMain}\essentiel\fontsize{11}{15}\selectfont
\phantomsection\label{cours-20-1}
\essSection{1}{Les tableaux}
\begin{definition}Un tableau range des données pour les retrouver vite. Dans un tableau à double entrée, une case se lit au croisement d’une ligne et d’une colonne : elle répond à deux questions à la fois.\end{definition}
\begin{exemple}Une enquête commence par une question (« Quel est ton fruit préféré ? »), continue par un recueil des réponses, puis par un comptage des effectifs : le tableau présente le résultat.\end{exemple}\begin{center}\renewcommand{\arraystretch}{1.35}\begin{tabular}{c|cccc}\toprule
Fruit & pomme & fraise & cerise & banane\\\midrule
Effectif & 7 & 9 & 5 & 4\\
\bottomrule\end{tabular}\end{center}
\begin{savoirfaire}[Contrôler les totaux]Contrôler un tableau d’enquête : la somme des effectifs doit retrouver le nombre de réponses. Ici $7+9+5+4=25$ élèves interrogés.\end{savoirfaire}
\clearpage
\phantomsection\label{cours-20-2}
\essSection{2}{Diagrammes et courbes}
\begin{definition}Dans un diagramme en barres, la hauteur de chaque barre représente un effectif, lu sur l’axe gradué. Dans une courbe, chaque point donne la valeur d’une grandeur à un instant : la ligne montre l’évolution.\end{definition}
\label{fig-cours-20-2-cours-donnees}\figureGeo{cours-donnees}

\begin{savoirfaire}[Lire les axes d’abord]Avant de lire : repérer le titre, les grandeurs des deux axes et le pas des graduations. Un carreau ne vaut pas toujours 1 : il peut valoir 2, 5, 10 ou plus.\end{savoirfaire}
\begin{attention}Les données réelles, par exemple des températures ou des précipitations, se lisent avec prudence : la courbe relie des mesures, elle ne dit rien d’exact entre deux points mesurés.\end{attention}

\clearpage\rubrique{Découverte}
\begin{blocLivre}{ex-20-1}\begin{exercice}[1]{Lire un tableau}
Le tableau du cours donne les fruits préférés de la classe. Combien d’élèves préfèrent la fraise ? Combien d’élèves ont répondu en tout ?
\end{exercice}\end{blocLivre}
\begin{blocLivre}{ex-20-2}\begin{exercice}[1]{Double entrée}
Le tableau croise deux classes et deux menus.\begin{center}\renewcommand{\arraystretch}{1.35}\begin{tabular}{c|cc}\toprule
 & Menu A & Menu B\\\midrule
CM1 & 17 & 8\\
CM2 & 14 & 12\\
\bottomrule\end{tabular}\end{center}Combien d’élèves de CM2 prennent le menu A ? Combien d’élèves prennent le menu B ?
\end{exercice}\end{blocLivre}
\begin{blocLivre}{ex-20-3}\begin{exercice}[1]{Compléter les totaux}
Reprendre le tableau des menus et calculer : le nombre d’élèves de CM1, le nombre d’élèves de CM2 et le nombre total d’élèves.
\end{exercice}\end{blocLivre}
\begin{blocLivre}{ex-20-4}\begin{exercice}[1]{Le diagramme en barres}
Le diagramme donne le nombre de livres empruntés par jour. Lire le nombre de livres du mardi et celui du jeudi. Quel jour en compte le plus ?
\label{fig-ex-20-4}\figureGeo{don-4}
\end{exercice}\end{blocLivre}
\begin{blocLivre}{ex-20-5}\begin{exercice}[1]{La courbe de température}
La courbe donne la température d’une journée. Lire la température à 8 h et à 14 h. À quelle heure la température est-elle la plus haute ?
\label{fig-ex-20-5}\figureGeo{don-5}
\end{exercice}\end{blocLivre}

\rubrique{Entraînement}
\begin{blocLivre}{ex-20-6}\begin{exercice}[2]{Construire un diagramme}
Construire le diagramme en barres des fruits préférés du cours : pomme 7, fraise 9, cerise 5, banane 4. L’axe est gradué de 1 en 1.
\label{fig-ex-20-6}\figureGeo{don-6}
\end{exercice}\end{blocLivre}
\begin{blocLivre}{ex-20-7}\begin{exercice}[2]{Des points dans un repère}
Chaque point du repère donne la masse d’un chiot selon son âge en semaines. Lire la masse à 2 semaines et à 6 semaines.
\label{fig-ex-20-7}\figureGeo{don-7}
\end{exercice}\end{blocLivre}
\begin{blocLivre}{ex-20-8}\begin{exercice}[2]{Des questions en étapes}
Avec le tableau des menus de l’exercice 2 : combien d’élèves de plus en CM2 qu’en CM1 ont choisi le menu B ? Le menu A est-il plus choisi que le menu B ?
\end{exercice}\end{blocLivre}
\begin{blocLivre}{ex-20-9}\begin{exercice}[2]{L’année des températures}
La courbe donne la température moyenne de chaque mois d’une ville. Quel est le mois le plus chaud ? Le plus froid ? Calculer l’écart entre les deux.
\label{fig-ex-20-9}\figureGeo{don-9}
\end{exercice}\end{blocLivre}
\begin{blocLivre}{ex-20-10}\begin{exercice}[2]{La barre fausse}
Le tableau annonce : lundi 8, mardi 12, mercredi 6 tickets vendus. Une barre du diagramme ne correspond pas au tableau. Laquelle ? Corriger.
\label{fig-ex-20-10}\figureGeo{don-10}
\end{exercice}\end{blocLivre}

\rubrique{Problèmes}
\begin{blocLivre}{ex-20-11}\begin{exercice}[3]{L’enquête brute}
Les réponses brutes à la question « Comment viens-tu à l’école ? » sont : à pied, vélo, à pied, bus, vélo, à pied, bus, à pied, vélo, à pied.\begin{questions}
\item Construire le tableau des effectifs.
\item Quel est le moyen le plus utilisé ?
\end{questions}
\end{exercice}\end{blocLivre}
\begin{blocLivre}{ex-20-12}\begin{exercice}[3]{La bibliothèque}
Le tableau donne les prêts de la bibliothèque.\begin{center}\renewcommand{\arraystretch}{1.35}\begin{tabular}{c|cc}\toprule
 & Romans & Documentaires\\\midrule
Semaine 1 & 34 & 21\\
Semaine 2 & 28 & 25\\
\bottomrule\end{tabular}\end{center}\begin{questions}
\item Combien de livres ont été prêtés en semaine 1 ?
\item Les prêts de documentaires ont-ils augmenté ou diminué ?
\item Combien de romans ont été prêtés en deux semaines ?
\end{questions}
\end{exercice}\end{blocLivre}
\begin{blocLivre}{ex-20-13}\begin{exercice}[3]{La plante qui pousse}
La courbe donne la hauteur d’une plante semée en classe, mesurée chaque semaine. Lire la hauteur à la semaine 3, puis dire à partir de quelle semaine la plante dépasse 10 cm.
\label{fig-ex-20-13}\figureGeo{don-13}
\end{exercice}\end{blocLivre}
\begin{blocLivre}{ex-20-14}\begin{exercice}[3]{Deux villes, deux courbes}
Le graphique superpose les températures de Brestel et de Lyonnet sur quatre saisons. Quelle ville est la plus chaude en été ? En hiver ? Quelle ville a les températures les plus régulières ?
\label{fig-ex-20-14}\figureGeo{don-14}
\end{exercice}\end{blocLivre}
\begin{blocLivre}{ex-20-15}\begin{exercice}[3]{Les précipitations}
Le tableau donne les précipitations d’une ville, en millimètres.\begin{center}\renewcommand{\arraystretch}{1.35}\begin{tabular}{c|cccc}\toprule
Mois & sept. & oct. & nov. & déc.\\\midrule
Pluie (mm) & 45 & 78 & 96 & 60\\
\bottomrule\end{tabular}\end{center}\begin{questions}
\item Quel mois est le plus pluvieux ?
\item Combien de pluie est tombée en tout sur ces quatre mois ?
\item Novembre a-t-il reçu plus de pluie que septembre et décembre réunis ?
\end{questions}
\end{exercice}\end{blocLivre}

\clearpage
